\documentclass[11pt]{article}
\usepackage{fontspec,amsmath,amssymb,amsthm,hyperref}
\usepackage[margin=1in]{geometry}
\setlength{\emergencystretch}{3em}
\title{Coupling lifts and the squared transport infimum}
\author{LangevinLean mathematical exposition}
\date{}
\begin{document}\maketitle
\noindent Source module: \texttt{AverageLMC/SquaredTransport.lean}. Compile this proof document with XeLaTeX or LuaLaTeX. The accompanying HTML pairs the mathematical argument with the exact Lean statements and records the scope of the formalization.


The SGLD argument begins with any finite-cost coupling of the full-gradient output and the Gibbs target. A regular conditional distribution lifts this coupling to the original initialization and Gaussian space. No optimal coupling is assumed.

\[W_2(\mu,\pi)<\infty,\quad r\le a\,C(\gamma)^2+b\quad\hbox{for every finite-cost coupling }\gamma\quad\Longrightarrow\quad r\le aW_2(\mu,\pi)^2+b.\]
\section{Pass real quadratic estimates through the infimum}


For a nonnegative coefficient, choose finite couplings with costs arbitrarily close to the Wasserstein infimum and pass to the limit by continuity. For a negative coefficient, any finite coupling suffices: multiplication reverses the comparison between its squared cost and the infimum. The transport specialization retains an explicit finiteness condition for the output before its real value is used.



\noindent Lean declarations: \texttt{real\_bound\_of\_finite\_couplings}, \texttt{wasserstein2\_sq\_bound\_of\_finite\_couplings}.

\section{Disintegrate the output coupling over its first marginal}


The conditional kernel of the chosen coupling is composed with the original output map. Integrating this kernel against the original input law constructs a probability measure on input and target. Its input marginal is the original law. Mapping input and target to output and target recovers the chosen coupling, by the disintegration identity; taking its second marginal therefore recovers the Gibbs target. These identities justify adding an independent stochastic-gradient index to an actual output coupling.



\noindent Lean declarations: \texttt{liftedOutputCoupling}, \texttt{liftedOutputCoupling\_probability}, \texttt{liftedOutputCoupling\_fst}, \texttt{liftedOutputCoupling\_output}, \texttt{liftedOutputCoupling\_snd}.

\end{document}
